% GENERATED FILE. Edit canonical lessons, then run npm run generate:books.
% canonical-source-hash: fnv1a64:0159d63a2c4c2da5
% canonical-lessons: GU-C143-bhagvu, GU-C143-pharvu, GU-C143-atakvu, GU-C143-manvu, GU-C143-badalvu, GU-R143-first-pass-describing-words, GU-R143-second-pass-verbs, GU-W10-a1-timed-production

\chapter{To run away, To go around, To stop, To believe, To change, A1 timed writing}
\label{ch:gu-to-run-away-to-go-around-to-stop-to-believe-to-change}
\hlchaptermodality{143}

\begin{chapteropening}
\textbf{By the end of this chapter:} \emph{I can say to run away, to go around, to stop, to believe and to change, then complete the A1 writing paper within 20 minutes.}

Complete the last lesson of chapter 143: I can fill a six-field form and write a 30–40 word message within the A1 paper's 20-minute limit.
\end{chapteropening}

\section[\texorpdfstring{\gu{ભાગવું} (bhāgvũ)}{ભાગવું (bhāgvũ)}]{\gu{ભાગવું} (bhāgvũ) — to run away}

\label{lesson:GU-C143-bhagvu}



\begin{quote}
Before the new one: say the Gujarati for to save, then the Gujarati for to sow.
\end{quote}

\subsection*{You'll want to know: \gu{ભાગવું}}
\textbf{\gu{ભાગવું}} (\emph{bhāgvũ}) — \textquotedblleft{}to run away\textquotedblright{}.

\begin{grammarlens}[title={What you've built}]
One more word for this chapter.
\end{grammarlens}

\subsection*{Your turn}
\begin{itemize}
\raggedright
  \item \emph{Say it:} \emph{bhāgvũ}
  \item \emph{Say it:} \emph{bhāgvũ}, once more
  \item \emph{Say it:} say \emph{bhāgvũ}
  \item \emph{Recall:} say the Gujarati for to save, then the Gujarati for to sow, then say \emph{bhāgvũ} again
\end{itemize}

\begin{tcolorbox}[breakable,colback=teal!4,colframe=teal!35!black,title={Before you move on}]
What does \gu{ભાગવું} mean? (\textquotedblleft{}To run away\textquotedblright{}.) Say it once more.
\end{tcolorbox}
\section[\texorpdfstring{\gu{ફરવું} (pharvũ)}{ફરવું (pharvũ)}]{\gu{ફરવું} (pharvũ) — to go around}

\label{lesson:GU-C143-pharvu}



\begin{quote}
Before the new one: say the Gujarati for to sow, then the Gujarati for to run away.
\end{quote}

\subsection*{You'll want to know: \gu{ફરવું}}
\textbf{\gu{ફરવું}} (\emph{pharvũ}) — \textquotedblleft{}to go around\textquotedblright{}.

\begin{grammarlens}[title={What you've built}]
One more word for this chapter.
\end{grammarlens}

\subsection*{Your turn}
\begin{itemize}
\raggedright
  \item \emph{Say it:} \emph{pharvũ}
  \item \emph{Say it:} \emph{pharvũ}, once more
  \item \emph{Say it:} say \emph{pharvũ}
  \item \emph{Recall:} say the Gujarati for to sow, then the Gujarati for to run away, then say \emph{pharvũ} again
\end{itemize}

\begin{tcolorbox}[breakable,colback=teal!4,colframe=teal!35!black,title={Before you move on}]
What does \gu{ફરવું} mean? (\textquotedblleft{}To go around\textquotedblright{}.) Say it once more.
\end{tcolorbox}
\section[\texorpdfstring{\gu{અટકવું} (aṭakvũ)}{અટકવું (aṭakvũ)}]{\gu{અટકવું} (aṭakvũ) — to stop}

\label{lesson:GU-C143-atakvu}



\begin{quote}
Before the new one: say the Gujarati for to run away, then the Gujarati for to go around.
\end{quote}

\subsection*{You'll want to know: \gu{અટકવું}}
\textbf{\gu{અટકવું}} (\emph{aṭakvũ}) — \textquotedblleft{}to stop\textquotedblright{}.

\begin{grammarlens}[title={What you've built}]
One more word for this chapter.
\end{grammarlens}

\subsection*{Your turn}
\begin{itemize}
\raggedright
  \item \emph{Say it:} \emph{aṭakvũ}
  \item \emph{Say it:} \emph{aṭakvũ}, once more
  \item \emph{Say it:} say \emph{aṭakvũ}
  \item \emph{Recall:} say the Gujarati for to run away, then the Gujarati for to go around, then say \emph{aṭakvũ} again
\end{itemize}

\begin{tcolorbox}[breakable,colback=teal!4,colframe=teal!35!black,title={Before you move on}]
What does \gu{અટકવું} mean? (\textquotedblleft{}To stop\textquotedblright{}.) Say it once more.
\end{tcolorbox}
\section[\texorpdfstring{\gu{માનવું} (mānvũ)}{માનવું (mānvũ)}]{\gu{માનવું} (mānvũ) — to believe}

\label{lesson:GU-C143-manvu}



\begin{quote}
Before the new one: say the Gujarati for to go around, then the Gujarati for to stop.
\end{quote}

\subsection*{You'll want to know: \gu{માનવું}}
\textbf{\gu{માનવું}} (\emph{mānvũ}) — \textquotedblleft{}to believe\textquotedblright{}.

\begin{grammarlens}[title={What you've built}]
One more word for this chapter.
\end{grammarlens}

\subsection*{Your turn}
\begin{itemize}
\raggedright
  \item \emph{Say it:} \emph{mānvũ}
  \item \emph{Say it:} \emph{mānvũ}, once more
  \item \emph{Say it:} say \emph{mānvũ}
  \item \emph{Recall:} say the Gujarati for to go around, then the Gujarati for to stop, then say \emph{mānvũ} again
\end{itemize}

\begin{tcolorbox}[breakable,colback=teal!4,colframe=teal!35!black,title={Before you move on}]
What does \gu{માનવું} mean? (\textquotedblleft{}To believe\textquotedblright{}.) Say it once more.
\end{tcolorbox}
\section[\texorpdfstring{\gu{બદલવું} (badalvũ)}{બદલવું (badalvũ)}]{\gu{બદલવું} (badalvũ) — to change}

\label{lesson:GU-C143-badalvu}



\begin{quote}
Before the new one: say the Gujarati for to stop, then the Gujarati for to believe.
\end{quote}

\subsection*{You'll want to know: \gu{બદલવું}}
\textbf{\gu{બદલવું}} (\emph{badalvũ}) — \textquotedblleft{}to change\textquotedblright{}.

\begin{grammarlens}[title={What you've built}]
One more word for this chapter.
\end{grammarlens}

\subsection*{Your turn}
\begin{itemize}
\raggedright
  \item \emph{Say it:} \emph{badalvũ}
  \item \emph{Say it:} \emph{badalvũ}, once more
  \item \emph{Say it:} say \emph{badalvũ}
  \item \emph{Recall:} say the Gujarati for to stop, then the Gujarati for to believe, then say \emph{badalvũ} again
\end{itemize}

\begin{tcolorbox}[breakable,colback=teal!4,colframe=teal!35!black,title={Before you move on}]
What does \gu{બદલવું} mean? (\textquotedblleft{}To change\textquotedblright{}.) Say it once more.
\end{tcolorbox}
\section[(dialogue)]{Describing words — a first pass}

\label{lesson:GU-R143-first-pass-describing-words}



\begin{quote}
Say the words for to believe and to change.
\end{quote}

\subsection*{Your turn}
Say these aloud:

\begin{itemize}
\raggedright
  \item empty, full, true, false, rich
  \item poor, young, strong, weak, necessary
  \item to jump, to swim, to climb, to open, to call
  \item to send, to forget, to look for, to use, to make
\end{itemize}

\begin{tcolorbox}[breakable,colback=teal!4,colframe=teal!35!black,title={Before you move on}]
Empty? (\textbf{\gu{ખાલી}}.) Full? (\textbf{\gu{ભરેલું}}.) True? (\textbf{\gu{સાચું}}.) False? (\textbf{\gu{ખોટું}}.) Rich? (\textbf{\gu{અમીર}}.) Poor? (\textbf{\gu{ગરીબ}}.) Young? (\textbf{\gu{જુવાન}}.) Strong? (\textbf{\gu{મજબૂત}}.) Weak? (\textbf{\gu{નબળું}}.) Necessary? (\textbf{\gu{જરૂરી}}.) To jump? (\textbf{\gu{કૂદવું}}.) To swim? (\textbf{\gu{તરવું}}.) To climb? (\textbf{\gu{ચઢવું}}.) To open? (\textbf{\gu{ખોલવું}}.) To call? (\textbf{\gu{બોલાવવું}}.) To send? (\textbf{\gu{મોકલવું}}.) To forget? (\textbf{\gu{ભૂલવું}}.) To look for? (\textbf{\gu{શોધવું}}.) To use? (\textbf{\gu{વાપરવું}}.) To make? (\textbf{\gu{બનાવવું}}.)
\end{tcolorbox}
\section[(dialogue)]{Verbs — a second pass}

\label{lesson:GU-R143-second-pass-verbs}



\begin{quote}
Say the words for to cut and to sew.
\end{quote}

\subsection*{Your turn}
Say these aloud:

\begin{itemize}
\raggedright
  \item to cut, to sew, to fill, to pull, to throw
  \item to wake up, to bathe, to wear, to drive, to count
  \item to say, to live, to spend, to save, to sow
  \item to run away, to go around, to stop, to believe, to change
\end{itemize}

\begin{tcolorbox}[breakable,colback=teal!4,colframe=teal!35!black,title={Before you move on}]
To cut? (\textbf{\gu{કાપવું}}.) To sew? (\textbf{\gu{સીવવું}}.) To fill? (\textbf{\gu{ભરવું}}.) To pull? (\textbf{\gu{ખેંચવું}}.) To throw? (\textbf{\gu{ફેંકવું}}.) To wake up? (\textbf{\gu{જાગવું}}.) To bathe? (\textbf{\gu{નાહવું}}.) To wear? (\textbf{\gu{પહેરવું}}.) To drive? (\textbf{\gu{ચલાવવું}}.) To count? (\textbf{\gu{ગણવું}}.) To say? (\textbf{\gu{કહેવું}}.) To live? (\textbf{\gu{જીવવું}}.) To spend? (\textbf{\gu{ખર્ચવું}}.) To save? (\textbf{\gu{બચાવવું}}.) To sow? (\textbf{\gu{વાવવું}}.) To run away? (\textbf{\gu{ભાગવું}}.) To go around? (\textbf{\gu{ફરવું}}.) To stop? (\textbf{\gu{અટકવું}}.) To believe? (\textbf{\gu{માનવું}}.) To change? (\textbf{\gu{બદલવું}}.)
\end{tcolorbox}
\section[(timed A1 writing)]{Your first complete A1 writing paper}

\label{lesson:GU-W10-a1-timed-production}



\begin{quote}
This is an attempt, not a model. You will write the answers, and this page will not show you any copyable answer.

Put away romanization, dictionaries, translators, and spell-checkers. Set one timer for \textbf{20 minutes}. Keep it running through both tasks.
\end{quote}

\subsection*{Writing — timed assessment production}
\#\#\# Task 1 — practical form: 7 minutes, 40 points

Complete \textbf{all six fields in Gujarati script} for a new neighbourhood activity group. Short answers are enough.

\begin{enumerate}
\raggedright
  \item Your name
  \item Your city or village
  \item A language you speak
  \item One place you like
  \item One drink or food you like
  \item One everyday action or kind of work you do
\end{enumerate}

At seven minutes, move on even if a field is unfinished. Do not restart the timer and do not polish the form with time meant for the message.

\#\#\# Task 2 — message: 13 minutes, 60 points

\textbf{Reader:} Mira. \textbf{Purpose:} help Mira recognise you and arrange your first meeting tomorrow.

Write \textbf{30–40 words in Gujarati script}. Greet Mira, introduce yourself, say where you live and which language you speak, mention one place and one drink or food you like, ask where Mira lives, and close politely. The instructions give content, not sentences: choose and connect the Gujarati yourself.

When the 20-minute timer rings, \textbf{stop}, even if a line is unfinished. Record the number of completed form fields and the message word count. An honest timed attempt is more useful than a perfect untimed one.

\begin{tcolorbox}[breakable,colback=teal!4,colframe=teal!35!black,title={Before you move on}]
Only after you have stopped, review the same attempt. Do not rewrite it yet. Mark one strength and one next step under each assessment criterion:

\begin{enumerate}
\raggedright
  \item \textbf{Task fulfilment and relevant content} — six filled fields; the named reader, meeting purpose, and requested message points are present.
  \item \textbf{Comprehensibility and organisation} — Mira can follow the message and see how its ideas connect.
  \item \textbf{Gujarati vocabulary and grammar} — familiar words and sentence patterns carry the intended meaning.
  \item \textbf{Orthographic control} — Gujarati letters, vowel signs, spacing, and punctuation remain readable.
\end{enumerate}

Keep the paper. On another day, repeat the same two-task shape with new personal details. The goal is not to memorise this prompt; it is to make a complete Gujarati attempt calmly inside the real limit.
\end{tcolorbox}
